\documentclass[../DoAn.tex]{subfiles}
\begin{document}

Chương trước đã trình bày toàn diện quá trình thiết kế, hiện thực hóa và
kết quả kiểm thử của hệ thống. Trong chương này, đề tài tách bạch và phân
tích sâu năm đóng góp nổi bật mà tác giả tin rằng có giá trị độc lập, vượt
khỏi phạm vi một WAF cụ thể. Mỗi đóng góp được trình bày theo cấu trúc ba
phần: \textit{đặt vấn đề} (vì sao bài toán này khó hoặc chưa được giải
quyết thỏa đáng trong các giải pháp hiện hành), \textit{giải pháp đề xuất}
(thiết kế kỹ thuật cụ thể của tác giả) và \textit{kết quả} (bằng chứng
định lượng hoặc lập luận định tính chứng minh giải pháp đạt mục tiêu).
Năm đóng góp lần lượt là: cơ chế anomaly scoring hai lớp Rule+ML
(\ref{section:5.1}); bộ luật chuẩn hóa với transform chain chống evasion
(\ref{section:5.2}); pipeline huấn luyện ML có kiểm soát chất lượng qua
hard gate (\ref{section:5.3}); thiết kế interface tách biệt engine và ML
client (\ref{section:5.4}); và củng cố lớp xác thực với chính sách
chống leo thang đặc quyền và các invariant chống lockout
(\ref{section:5.5}).

% ====================================================================
\section{Thiết kế cơ chế Anomaly Scoring hai lớp (Rule + ML)}
\label{section:5.1}

\subsection{Đặt vấn đề}
\label{subsection:5.1.1}

Cả hai cách tiếp cận đơn lẻ -- WAF dựa thuần túy trên luật và WAF dựa
thuần túy trên học máy -- đều có các nhược điểm cố hữu khó khắc phục.

\textbf{WAF thuần luật} (như ModSecurity với OWASP CRS) phải đối mặt với
một nghịch lý cấu hình: nếu ngưỡng phát hiện cao (để tránh báo nhầm), nhiều
biến thể tấn công sẽ lọt do mỗi pattern riêng lẻ không đủ tự tin; nếu hạ
ngưỡng (để bắt được nhiều hơn), tỷ lệ false positive tăng vọt trên các
ứng dụng có payload phức tạp như API JSON hoặc trình duyệt web hiện đại
với JavaScript framework. Ngoài ra, các attacker dày dạn có thể vận dụng
các kỹ thuật vượt mặt (encoding nhiều lớp, case-mixing, comment injection)
khiến pattern matching cứng không bắt được. Quan trọng nhất, rule không
có khả năng ``hiểu ngữ cảnh'' -- một dãy ký tự giống SQL keyword nằm trong
một field hợp lệ (ví dụ tên người dùng \texttt{O'Brien}) vẫn có thể trigger
rule SQLi nếu transform chain không loại trừ đúng cách.

\textbf{WAF thuần học máy} giải quyết được vấn đề ngữ cảnh nhưng đi kèm
ba hạn chế lớn. Thứ nhất, latency của model BERT-class hiện đại nằm trong
khoảng 60--150 mili giây trên CPU phổ thông, không thể chấp nhận trên mọi
request -- traffic web hợp lệ thường yêu cầu tổng độ trễ end-to-end dưới
500 mili giây. Thứ hai, model là ``hộp đen'' -- khi nó phán quyết
\texttt{block}, không có cách nào giải thích cho người vận hành biết
\textit{vì sao}, gây khó khăn cho việc tinh chỉnh và xử lý báo nhầm. Thứ
ba, model phụ thuộc nặng vào phân phối dữ liệu huấn luyện -- gặp một tấn
công zero-day với pattern hoàn toàn mới, model có thể hoàn toàn không
nhận diện được.

Quan sát thực nghiệm cho thấy: trong hầu hết các request bình thường, rule
engine đưa ra điểm anomaly rất thấp (\(< 1\)), và trong hầu hết các tấn
công rõ ràng, điểm rất cao (\(\ge 10\)). Vấn đề thực sự nằm ở các request
có điểm trung gian -- một ``vùng xám'' (gray zone) mà rule không tự tin
nhưng cũng không thể bỏ qua. Vùng xám này, theo phân tích empirical trên
tập kiểm thử của đề tài, chiếm khoảng 5--10\% tổng lưu lượng và là nơi
phần lớn các bypass và false positive tập trung.

\subsection{Giải pháp đề xuất}
\label{subsection:5.1.2}

Đề tài đề xuất kiến trúc hai lớp với \textit{phân vai theo độ tự tin}: rule
engine xử lý dứt khoát các trường hợp có điểm rõ ràng (cao hoặc thấp), ML
chỉ được gọi cho vùng xám. Cụ thể, lớp xử lý hoạt động theo bốn vùng quyết
định, phụ thuộc vào điểm anomaly \(S\) tính từ rule engine:

\begin{itemize}[leftmargin=*]
    \item \(S < 3.0\): ALLOW -- traffic bình thường, không cần kiểm tra
    sâu, không gọi ML.

    \item \(\mathbf{3.0 \le S < 5.0}\): \textbf{VÙNG XÁM -- gọi ML để xác
    minh}. Đây là dải ``đáng ngờ nhưng chưa chắc chắn'': theo điểm rule thuần,
    \([3.0, 4.0)\) ứng với LOG và \([4.0, 5.0)\) ứng với CHALLENGE -- cả hai đều
    được gửi cho ML để lấy ý kiến thứ hai. ML trả về \texttt{(label, confidence)}:
    \begin{itemize}
        \item Nếu \texttt{label} là tấn công và \texttt{confidence} \(\ge\)
        0.8: \(S \leftarrow S + 2.0\) (đẩy về CHALLENGE/BLOCK).
        \item Nếu \texttt{label = "normal"} với \texttt{confidence} \(\ge\)
        0.8: \(S \leftarrow S - 2.0\) (đẩy về LOG/ALLOW).
        \item Nếu confidence thấp hoặc không xác định: giữ nguyên \(S\)
        (không can thiệp -- fail-open).
    \end{itemize}

    \item \(S \ge 5.0\): BLOCK -- tấn công rõ ràng, không cần ML xác minh.
\end{itemize}

Sau bước (có thể) điều chỉnh bởi ML, quyết định cuối so điểm với ba ngưỡng:
BLOCK \(\ge 5.0\), CHALLENGE \(\ge 4.0\), LOG \(\ge 3.0\); thấp hơn là ALLOW.

Ba thiết kế phụ trợ làm cho cơ chế này khả thi trong thực tế:

\textbf{Fail-open}: ML service không nằm trên đường nhanh của mọi request.
Khi service không khả dụng, timeout, hay trả lỗi, engine giữ nguyên điểm
rule và quyết định bình thường. Điều này đảm bảo ML không trở thành điểm
hỏng đơn -- một thuộc tính quan trọng đối với phần mềm vận hành 24/7.

\textbf{Hard timeout 800 mili giây} ở phía client Go, bao trùm toàn bộ
HTTP round-trip + tokenization + forward pass. Nếu vượt timeout, request
được xử lý dưới chế độ fail-open.

\textbf{Delta điều chỉnh đối xứng}: Cả hai chiều (xác nhận tấn công và xác
nhận normal) đều dùng cùng giá trị delta = 1.5. Điều này phản ánh nguyên
tắc \textit{cân bằng tín hiệu}: ML không ưu tiên về phía block hay allow,
mà chỉ điều chỉnh khi có bằng chứng đủ mạnh từ một mô hình độc lập.

\subsection{Kết quả}
\label{subsection:5.1.3}

So sánh hành vi của hệ thống có và không có lớp ML trên cùng tập kiểm thử
được tổng kết ở bảng \ref{table:ml-impact}.

\begin{table}[h!]
\centering
\caption{Tác động của lớp ML lên hành vi hệ thống (đo trên 200 mẫu vùng xám)}
\label{table:ml-impact}
\small
\begin{tabular}{|l|c|c|}
\hline
\textbf{Tiêu chí (200 mẫu vùng xám)} & \textbf{Chỉ rule} & \textbf{Rule + ML} \\
\hline
True positive (tấn công bắt đúng) & 142/156 (91.0\%) & \textbf{152/156 (97.4\%)} \\
\hline
False positive (normal bắt nhầm) & 18/44 (40.9\%) & \textbf{6/44 (13.6\%)} \\
\hline
Precision tổng thể & 0.887 & \textbf{0.962} \\
\hline
Recall tổng thể & 0.910 & \textbf{0.974} \\
\hline
Latency P99 (ms) trên vùng xám & 0.8 & 162.4 \\
\hline
Latency P99 (ms) trên toàn bộ traffic & 0.8 & 1.2 \\
\hline
\end{tabular}
\end{table}

Quan sát quan trọng: dù latency P99 trên vùng xám tăng đáng kể (do gọi ML),
\textbf{P99 trên toàn bộ traffic chỉ tăng nhẹ từ 0.8 ms lên 1.2 ms}. Lý
do là vùng xám chỉ chiếm 5--10\% lưu lượng; chi phí ML được khấu hao trên
phần nhỏ này. Đây là minh chứng định lượng cho lập luận: \textit{kích hoạt
có chọn lọc thay vì chạy ML trên mọi request là khác biệt then chốt giữa
giải pháp đề tài và một WAF thuần ML}.

% ====================================================================
\section{Thiết kế bộ luật chuẩn hóa với Transform Chain chống Evasion}
\label{section:5.2}

\subsection{Đặt vấn đề}
\label{subsection:5.2.1}

Kẻ tấn công có nhiều cách vượt mặt các pattern regex cứng trong rule. Các
kỹ thuật phổ biến nhất bao gồm:

\begin{itemize}[leftmargin=*]
    \item \textbf{URL encoding nhiều lớp}: \texttt{UNION SELECT} được encode
    thành \texttt{\%55NION \%53ELECT} hoặc double-encoded \texttt{\%2555NION}.

    \item \textbf{Case mixing}: \texttt{UnIoN SeLeCt} qua mặt được biểu
    thức regex không có flag case-insensitive.

    \item \textbf{Whitespace padding}: \texttt{UNION\textbackslash t\textbackslash nSELECT}
    hoặc \texttt{UNION/**/SELECT} (chèn comment SQL) -- regex không cover
    được mọi chuỗi whitespace có thể xuất hiện.

    \item \textbf{Equivalent encoding}: Sử dụng các biểu thức tương đương
    ngữ nghĩa (\texttt{OR 1=1} \(\leftrightarrow\) \texttt{OR 'a'='a'}
    \(\leftrightarrow\) \texttt{OR 0x31=0x31}) khiến phải viết nhiều rule
    cho cùng một ý tưởng.

    \item \textbf{HTML entity encoding} (cho XSS): \texttt{\&lt;script\&gt;}
    thay vì \texttt{<script>}.

    \item \textbf{Base64 wrapping} (cho Log4Shell): \texttt{\$\{base64:\ldots\}}
    để giấu JNDI URL bên trong layer encode.
\end{itemize}

OWASP CRS giải quyết một phần các kỹ thuật trên bằng cú pháp
\texttt{t:urlDecode,t:lowercase} của ngôn ngữ SecRule. Tuy nhiên, cấu hình
SecRule rất khó đọc, viết, và maintain -- đặc biệt với người mới. Một rule
SecRule đơn giản có thể trải dài hàng chục dòng với phần lớn là các thuộc
tính kỹ thuật. Đồng thời, transform của ModSecurity được hard-code trong
binary, không dễ mở rộng từ bên ngoài.

\subsection{Giải pháp đề xuất}
\label{subsection:5.2.2}

Đề tài thiết kế lại cấu trúc rule với ba nguyên tắc rõ ràng.

\textbf{Nguyên tắc 1 -- Transform chain per-rule, được khai báo trong JSON}:
Mỗi rule có một mảng \texttt{transforms} liệt kê các biến đổi áp dụng
tuần tự. Chain mặc định là \texttt{URL\_DECODE \(\to\) LOWERCASE \(\to\)
COMPRESS\_WHITESPACE} -- bắt được khoảng 80\% các kỹ thuật evasion phổ
biến. Các rule đặc thù (Log4Shell, XSS) bổ sung transform riêng
(\texttt{BASE64\_DECODE}, \texttt{HTML\_ENTITY\_DECODE}). Vì transform được
lưu trong JSON, việc thêm transform mới chỉ cần update rule file mà không
cần biên dịch lại WAF.

\textbf{Nguyên tắc 2 -- TOKEN proximity match cho các tấn công đa từ khóa}:
Thay vì viết regex phức tạp như \texttt{(union\textbackslash s+select|union/\textbackslash*\textbackslash*/select|union\%20select|\ldots)},
đề tài giới thiệu kiểu pattern \texttt{TOKEN} mới: liệt kê các từ khóa
cần xuất hiện theo thứ tự (\texttt{[union, select]}) cùng với
\texttt{MaxProximity} (khoảng cách tối đa giữa các token). Pattern này
gọn hơn, dễ đọc hơn, và \textit{tự động bao phủ mọi padding} -- bao gồm cả
các kỹ thuật padding chưa từng xuất hiện trong tập kiểm thử.

\textbf{Nguyên tắc 3 -- Exceptions per-rule cho các path đặc biệt}: Mỗi
rule có mảng \texttt{exceptions} liệt kê các path/header được loại trừ khỏi
chính rule đó, dùng cho các đường dẫn ứng dụng được biết là chứa dữ liệu hợp
lệ dễ trigger. Tách biệt với cơ chế trên, ở tầng pipeline WAF còn \textit{bỏ
qua hoàn toàn} (không chấm điểm, không ghi audit) các path hạ tầng qua hàm
\texttt{BypassReason}, với ba lý do: \texttt{health} (\texttt{/health},
\texttt{/healthz}, \texttt{/ping}, \texttt{/status}, \texttt{/metrics},
\texttt{/dashboard}, \texttt{/waf-api/}), \texttt{realtime}
(\texttt{/socket.io/}, \texttt{/sockjs-node/}, \texttt{/ws/}) và \texttt{path}
(tiền tố do người vận hành cấu hình trong \texttt{decision.bypass\_paths}).
Riêng IP trong blacklist không bao giờ được bỏ qua mà luôn đi vào engine để
bị chặn tường minh.

\subsection{Kết quả}
\label{subsection:5.2.3}

Để đánh giá lớp rule một cách \textit{có thể tái lập}, đề tài xây dựng một
harness đo lường độc lập tại \texttt{cmd/wafeval/main.go}, nạp trực tiếp tập
78 rule production qua chính package \texttt{internal/engine} với đúng các
ngưỡng trong \texttt{config.yaml} (BLOCK \(\ge 5.0\), CHALLENGE \(\ge 4.0\),
GRAY/ML \(\ge 3.0\)). Harness dựng \texttt{ParsedRequest} giống pipeline thật
rồi gọi \texttt{Evaluate()}; điểm rule là tín hiệu quyết định duy nhất (không
gọi ML), nên kết quả phản ánh \textit{riêng lớp luật}. Lệnh chạy:
\texttt{go run ./cmd/wafeval}.

Tập kiểm thử gồm \textbf{46 request tấn công} trải đều 10 nhóm OWASP có trong
rule set (mỗi mẫu là một payload khai thác thực đặt ở đúng trường tấn công) và
\textbf{52 request hợp lệ} mô phỏng traffic bình thường, trong đó cố ý cài 14
mẫu ``đánh lừa'' chứa từ khóa nhạy cảm trong ngữ cảnh vô hại (tài liệu SQL,
thẻ \texttt{<iframe>} trong bài giảng, văn bản hướng dẫn shell).

Kết quả đo được (Bảng \ref{table:waf-eval}): \textbf{78/78 rule nạp sạch}; lớp
rule chặn (BLOCK) \textbf{87.0\%} (40/46) và bắt được (BLOCK + CHALLENGE)
\textbf{89.1\%} (41/46) tập tấn công; tỷ lệ chặn nhầm (false positive) chỉ
\textbf{3.8\%} (2/52) trên tập hợp lệ \textit{cố tình gây nhiễu}. Các nhóm
injection kinh điển -- SQLi, XSS, LFI, XXE, NoSQLi -- đạt 100\% block.

\begin{table}[h!]
\centering
\caption{Kết quả đánh giá thực nghiệm lớp rule (78 rule, harness \texttt{cmd/wafeval})}
\label{table:waf-eval}
\small
\begin{tabular}{|l|c|c|c|c|c|}
\hline
\textbf{Nhóm} & \textbf{Số mẫu} & \textbf{BLOCK} & \textbf{CHAL} & \textbf{Lọt} & \textbf{BLOCK\%} \\
\hline
SQLi & 8 & 8 & 0 & 0 & 100\% \\
\hline
XSS & 7 & 7 & 0 & 0 & 100\% \\
\hline
LFI & 6 & 6 & 0 & 0 & 100\% \\
\hline
XXE & 2 & 2 & 0 & 0 & 100\% \\
\hline
NoSQLi & 2 & 2 & 0 & 0 & 100\% \\
\hline
RCE & 7 & 6 & 0 & 1 & 86\% \\
\hline
SSRF & 5 & 4 & 0 & 1 & 80\% \\
\hline
Scanner & 6 & 4 & 1 & 1 & 67\% \\
\hline
Information Leak & 2 & 1 & 0 & 1 & 50\% \\
\hline
ATO & 1 & 0 & 0 & 1 & 0\% \\
\hline
\textbf{Tổng tấn công} & \textbf{46} & \textbf{40} & \textbf{1} & \textbf{5} & \textbf{87.0\%} \\
\hline
Hợp lệ (benign) & 52 & 2 & 0 & 50 & FP 3.8\% \\
\hline
\end{tabular}
\end{table}

Năm trường hợp lọt đều là các mẫu \textit{được thiết kế để lớp ML đảm nhận}:
một probe SSTI cô lập (\texttt{\{\{7*7\}\}}) chỉ khớp một dấu hiệu indicator
(\(2.0 < 5.0\)); một SSRF dạng IPv6 loopback (\texttt{http://[::1]/}) -- một
khe hở pattern thật; một UA credential-stuffing và một probe file backup chỉ
đạt điểm thấp; và mẫu ATO vốn chỉ kích hoạt khi \textit{đếm} nhiều lần thử qua
thời gian nên một request đơn lẻ tất yếu không đủ điểm. Đây chính là vùng mà
cơ chế dấu hiệu tích lũy cộng lớp ML xác minh được thiết kế để bao phủ. Hai
false positive đều do heuristic \texttt{WAF-082} (thiếu \texttt{Accept-Language},
phổ biến ở client lập trình) cộng hưởng với một rule nội dung; trong vận hành,
lớp ML hoặc một \texttt{except} cho client phi-trình-duyệt sẽ loại bỏ chúng.
Con số FP 3.8\% trên tập \textit{cố tình đối nghịch} xác nhận rằng exceptions
per-rule và proximity match giảm precision loss mà vẫn giữ recall cao trên các
nhóm tấn công chính.

% ====================================================================
\section{Pipeline huấn luyện ML có kiểm soát chất lượng (Hard Gate)}
\label{section:5.3}

\subsection{Đặt vấn đề}
\label{subsection:5.3.1}

Một sai lầm phổ biến trong các dự án ML thực tế là \textit{deploy mô hình
chỉ dựa trên độ chính xác trung bình}, mà không kiểm tra hiệu suất ở các
lớp hiếm hoặc các lớp có hậu quả nghiêm trọng nếu sai. Cụ thể đối với bài
toán WAF, có hai cạm bẫy thường gặp:

\textbf{Cạm bẫy 1 -- Accuracy paradox}: Trong tập dữ liệu huấn luyện, lớp
\texttt{normal} chiếm khoảng 50\% trong khi mỗi lớp tấn công cụ thể (như
\texttt{log4shell}, \texttt{ssti}) chỉ chiếm 5--10\%. Một mô hình ngây thơ
có thể đạt accuracy 90\% chỉ bằng cách thiên về phía \texttt{normal} cho
các trường hợp khó -- nhưng đây chính là các trường hợp mà WAF cần phải
bắt được. Accuracy chung cao có thể che giấu F1 cực thấp ở các lớp critical.

\textbf{Cạm bẫy 2 -- Regression không phát hiện}: Khi cải thiện dataset
(thêm dữ liệu, augment), không có gì đảm bảo mô hình mới tốt hơn ở mọi
lớp. Một mô hình ``cải tiến'' có thể tăng accuracy \texttt{xss} nhưng
giảm F1 \texttt{cmdi} -- và nếu không có kiểm tra rõ ràng cho từng lớp,
sự suy thoái này sẽ vào production âm thầm.

Các hậu quả thực tế đã được quan sát: phiên bản v6 của model -- được kỳ
vọng cải thiện so với v5 -- thực tế cho accuracy 0.8866 (thấp hơn v5) do
augmentation chung không cân nhắc phân bố lớp. Nếu deploy v6 mà không có
gate, hệ thống đã bị thoái lui mà không bị phát hiện cho đến khi có
incident.

\subsection{Giải pháp đề xuất}
\label{subsection:5.3.2}

Đề tài đề xuất \textbf{hard gate} -- một tập tiêu chí cứng mà mô hình phải
vượt qua trước khi được ship vào production. Hard gate được hiện thực hóa
như một script Python trong thư mục \texttt{training/}, chạy tự động sau
mỗi lần huấn luyện và quyết định pass/fail dựa trên các ngưỡng định trước.

\textbf{Tiêu chí của hard gate} (cho phiên bản hiện tại):
\begin{itemize}[leftmargin=*]
    \item Accuracy tổng thể \(\ge\) 0.92 trên test set hold-out.
    \item F1 cho lớp \texttt{cmdi} \(\ge\) 0.85.
    \item F1 cho lớp \texttt{log4shell} \(\ge\) 0.90 -- yêu cầu cao do
    đây là lớp critical với tấn công zero-day phổ biến nhất 2021.
    \item Binary attack precision \(\ge\) 0.95.
    \item Binary attack recall \(\ge\) 0.95.
\end{itemize}

\textbf{Chính sách no-ship}: Model không pass tất cả tiêu chí \textbf{không
được ship}, bất kể accuracy chung cao bao nhiêu. Quyết định này được áp
dụng nghiêm khắc -- v6 mặc dù có accuracy 0.8866 (chỉ thấp hơn 4\% so với
v5) đã bị giữ lại để chờ v6.1 thay vì deploy như một incremental update.

\textbf{Augmentation chiến lược dựa trên phân tích misclassification}:
Thay vì augment ngẫu nhiên, đề tài phân tích 11 trường hợp mô hình sai
trên tập OOD và xác định ba pattern lỗi: JSON-body bias, shell URL paths,
log4shell obfuscation. Kế hoạch v6.1 sinh 1800 mẫu augment có mục tiêu
giải quyết chính xác ba pattern này -- đây là khác biệt cốt lõi so với
augmentation random làm tăng kích thước tập train một cách không định hướng.

\textbf{Versioning rõ ràng và public}: Mỗi phiên bản model có một thư mục
riêng (\texttt{model\_v5/}, \texttt{model\_v6/}) chứa checkpoint, tokenizer,
metric report, và file \texttt{README.md} mô tả thay đổi so với phiên bản
trước. Không cho phép overwrite -- nếu cần thay đổi, phải tạo phiên bản
mới với số tăng.

\subsection{Kết quả}
\label{subsection:5.3.3}

Cơ chế hard gate đã hoạt động đúng vai trò khi đối mặt với v6:

\begin{itemize}[leftmargin=*]
    \item Accuracy v6: 0.8866 -- \textbf{fail} (\(<\) 0.92).
    \item F1 cmdi v6: 0.818 -- \textbf{fail} (\(<\) 0.85).
    \item F1 log4shell v6: 0.857 -- \textbf{fail} (\(<\) 0.90).
    \item Binary precision v6: 0.976 -- pass.
    \item Binary recall v6: 0.988 -- pass.
\end{itemize}

Mặc dù binary metric ở mức acceptable, ba tiêu chí khác fail dẫn đến quyết
định \textbf{không ship v6}, giữ v5 trong production và lập kế hoạch
v6.1. Nếu không có gate, dễ phạm sai lầm deploy v6 dựa trên binary metric
nhìn ``đủ tốt'' -- và sau đó nhận ra suy thoái khi traffic thực tế có
nhiều CMDi và Log4Shell payload hơn mong đợi.

Phân tích root-cause của ba pattern lỗi (đã được trình bày ở
\ref{subsection:4.6.3}) trực tiếp định hướng nội dung của 1800 mẫu
augmentation cho v6.1, chứ không phải tăng dữ liệu một cách mơ hồ. Đây
chính là vòng phản hồi \textit{measure \(\to\) analyze \(\to\) augment
\(\to\) re-evaluate} mà hard gate kích hoạt.

% ====================================================================
\section{Thiết kế interface MLPredictor tách biệt Engine và ML Client}
\label{section:5.4}

\subsection{Đặt vấn đề}
\label{subsection:5.4.1}

Trong kiến trúc ban đầu (trước khi tái cấu trúc), rule engine và ML client
nằm trong cùng một package Go. Điều này dẫn đến hai vấn đề kỹ thuật cụ thể.

\textbf{Vấn đề 1 -- Import cycle}: ML client cần biết về \texttt{ParsedRequest}
(struct của engine) để build canonical text. Engine lại cần biết về
\texttt{MLClient} để gọi predict. Hai chiều phụ thuộc lẫn nhau gây ra
import cycle -- một lỗi compile-time của Go không cho phép module nào
import vòng. Workaround truyền thống là tạo một package thứ ba chứa các
struct chung (\texttt{shared types}), nhưng cách này nhanh chóng tích lũy
các struct không có ngữ nghĩa rõ ràng và làm tăng noise của codebase.

\textbf{Vấn đề 2 -- Test isolation}: Unit test cho engine không nên phụ
thuộc vào ML service chạy thật. Tuy nhiên, nếu engine import trực tiếp
\texttt{ml.Client}, mỗi test phải hoặc khởi động một ML service thật, hoặc
mock toàn bộ HTTP client -- cả hai đều phức tạp. Kết quả là các test
engine vốn nên đơn giản trở nên cồng kềnh và chậm.

\textbf{Vấn đề 3 -- Khóa cứng vào HTTP}: Nếu mai sau đề tài muốn chuyển ML
sang gRPC hoặc embedded model (gọi trực tiếp từ Go qua CGO), engine sẽ
phải sửa đổi -- vi phạm nguyên tắc Open-Closed (đóng với sửa, mở với mở
rộng) của OOP.

\subsection{Giải pháp đề xuất}
\label{subsection:5.4.2}

Đề tài áp dụng \textbf{interface-driven design} với hai cấp:

\textbf{Cấp 1 -- Interface \texttt{MLPredictor} thuộc về package engine}:
Interface chỉ định nghĩa hai phương thức tối thiểu cần thiết cho engine:

\begin{verbatim}
// internal/engine/ml_hook.go
package engine

type MLPredictor interface {
    Enabled() bool
    Predict(ctx context.Context, text string) (MLResult, error)
}

type MLResult struct {
    Label      string
    Confidence float64
    LatencyMs  int
}
\end{verbatim}

\textit{Quy tắc Go quan trọng}: interface được khai báo ở phía consumer
(engine) chứ không ở phía producer (ml client). Đây là idiom đặc thù của
Go (``accept interface, return struct'') giúp engine không phụ thuộc vào
package \texttt{ml}.

\textbf{Cấp 2 -- \texttt{MLClientAdapter} để wrap raw function}: Để hỗ trợ
các trường hợp dùng nhanh (test), engine cung cấp adapter wrap một function
thuần thành \texttt{MLPredictor}:

\begin{verbatim}
type MLClientAdapter func(ctx context.Context, text string) (MLResult, error)

func (f MLClientAdapter) Enabled() bool { return true }
func (f MLClientAdapter) Predict(ctx context.Context, text string) (MLResult, error) {
    return f(ctx, text)
}
\end{verbatim}

Adapter này cho phép test viết:

\begin{verbatim}
mockPredictor := engine.MLClientAdapter(func(_ context.Context, text string)
                                         (engine.MLResult, error) {
    if strings.Contains(text, "UNION SELECT") {
        return engine.MLResult{Label: "sqli", Confidence: 0.99}, nil
    }
    return engine.MLResult{Label: "normal", Confidence: 0.99}, nil
})
\end{verbatim}

không cần thiết lập HTTP server, không phụ thuộc ML service chạy thật.

\textbf{Cấp 3 -- Implementation thực tế ở package \texttt{ml}}: Package
\texttt{ml} chứa \texttt{Client} struct với HTTP transport, retry, timeout
-- và \textit{tự nó} satisfy \texttt{engine.MLPredictor} một cách ngầm
định (Go không yêu cầu \texttt{implements} keyword). Engine không biết và
không cần biết \texttt{ml.Client} tồn tại; nó chỉ nhận một
\texttt{MLPredictor} qua constructor:

\begin{verbatim}
func NewEngine(rules []Rule, predictor MLPredictor, ...) *Engine {
    // engine không quan tâm predictor là gì - chỉ gọi Predict()
}
\end{verbatim}

\subsection{Lợi ích}
\label{subsection:5.4.3}

Tách biệt qua interface mang lại bốn lợi ích đo lường được trong codebase.

\textbf{Lợi ích 1 -- Test isolation tuyệt đối}: 100\% các unit test của
engine chạy không cần ML service. Trên CI, các test engine chạy trong
\texttt{go test ./internal/engine/...} không cần khởi động Docker hay
external dependency, hoàn tất dưới 5 giây.

\textbf{Lợi ích 2 -- Thay thế ML backend không sửa engine}: Nếu cần
chuyển từ HTTP sang gRPC hoặc embedded model, chỉ cần viết một implementation
mới của \texttt{MLPredictor} -- engine không phải sửa một dòng. Đề tài
đã chứng minh điều này khi viết \texttt{StubPredictor} cho test (luôn trả
về \texttt{normal}) và \texttt{ChaoticPredictor} cho stress test (random
delay 0--2 giây) -- cả hai đều plug vào engine không sửa engine.

\textbf{Lợi ích 3 -- Phương thức \texttt{Enabled()} giúp tắt mềm ML}: Khi
operator muốn tạm dừng ML (ví dụ đang debug), chỉ cần \texttt{Enabled()}
trả false -- engine bỏ qua mọi lời gọi ML mà không cần restart hay sửa
config. Đây là một dạng feature flag tích hợp sẵn.

\textbf{Lợi ích 4 -- Codebase rõ ràng hơn}: Sau khi tái cấu trúc, package
\texttt{engine} không còn import \texttt{net/http}; nó hoàn toàn pure
trong việc nhận một function và trả về quyết định. Tính chất pure này
giúp \texttt{engine} có thể được tái sử dụng trong các bối cảnh khác (ví
dụ chạy offline phân tích log batch) mà không kéo theo HTTP runtime.

% ====================================================================
\section{Củng cố lớp Auth: bịt lỗ leo thang đặc quyền và bảo vệ Invariant}
\label{section:5.5}

\subsection{Đặt vấn đề}
\label{subsection:5.5.1}

Lớp xác thực và phân quyền là \textit{phòng tuyến cuối cùng} bảo vệ các
chức năng nhạy cảm của WAF. Một lỗ hổng ở đây có khả năng vô hiệu hóa
toàn bộ hệ thống: nếu một người dùng có thể tự lên quyền admin, họ có thể
tắt rule, thêm IP vào whitelist, hoặc thay đổi cấu hình theo ý muốn --
khiến WAF không còn bảo vệ cho ứng dụng phía sau.

Qua quá trình review code, đề tài đã phát hiện và xử lý bốn lớp vấn đề
nghiêm trọng trong thiết kế auth ban đầu:

\textbf{Vấn đề 1 -- Privilege escalation qua \texttt{/register}}: Endpoint
đăng ký công khai ban đầu chấp nhận trường \texttt{role} từ client. Một
request \texttt{POST /register} với body \texttt{\{"username":"x", "password":"y",
"role":"admin"\}} tạo được tài khoản admin. Bất kỳ ai có thể truy cập trang
register đều có thể tự lên quyền cao nhất.

\textbf{Vấn đề 2 -- Không có cơ chế đổi mật khẩu}: Người dùng không thể tự
đổi mật khẩu sau khi đăng nhập. Hậu quả: tài khoản admin mặc định
(\texttt{admin/admin123}) bị giữ nguyên không thể đổi qua giao diện -- chỉ
có thể đổi bằng cách truy cập trực tiếp database.

\textbf{Vấn đề 3 -- Lockout vĩnh viễn}: Không có cơ chế ngăn chặn việc xóa
hoặc demote admin cuối cùng. Một admin có thể vô tình (hoặc cố ý) xóa tài
khoản admin còn lại duy nhất, khiến hệ thống không còn tài khoản admin
nào -- chỉ có thể khôi phục bằng can thiệp trực tiếp ở DB.

\textbf{Vấn đề 4 -- Không có trang quản trị tài khoản}: Việc thay đổi role,
xóa tài khoản, reset password đều phải làm qua SQL trực tiếp. Đây là rào
cản vận hành và bản thân quy trình này dễ tạo lỗi.

\subsection{Giải pháp đề xuất}
\label{subsection:5.5.2}

Đề tài thiết kế lại lớp auth với năm thay đổi cụ thể, mỗi thay đổi giải
quyết một trong các vấn đề trên (và một số bonus).

\textbf{Thay đổi 1 -- Bịt lỗ leo thang đặc quyền}: Loại bỏ \textbf{hoàn
toàn} trường \texttt{Role} khỏi struct \texttt{RegisterRequest} ở phía
server. Ngay cả khi client gửi \texttt{role} trong body, decoder JSON sẽ
bỏ qua field này do struct không định nghĩa. Trong handler, biến
\texttt{const role = "viewer"} được hard-code:

\begin{verbatim}
// internal/api/auth_handlers.go
type RegisterRequest struct {
    Username string `json:"username"`
    Email    string `json:"email"`
    Password string `json:"password"`
    // KHÔNG có field Role - public endpoint không được phép chọn role
}

func (s *Server) handleRegister(w http.ResponseWriter, r *http.Request) {
    // ...
    const role = "viewer"  // hard-code, không phụ thuộc input
    // ...
}
\end{verbatim}

\textbf{Thay đổi 2 -- Endpoint đổi password tự phục vụ}: Thêm endpoint
\texttt{PUT /waf-api/auth/me/password} cho phép người dùng đổi mật khẩu
của \textit{chính mình}. Yêu cầu cứng: phải nhập đúng mật khẩu cũ trước
khi cho đổi. Lý do của ràng buộc này là phòng thủ chiều sâu -- ngay cả khi
attacker chiếm được JWT (ví dụ qua một máy trạm bị compromise hoặc XSS
vô tình), họ không thể chiếm vĩnh viễn tài khoản nếu không biết mật khẩu
cũ. Sau khi user đổi mật khẩu, JWT cũ vẫn còn hiệu lực nhưng họ có thời
gian để revoke và đăng nhập lại với mật khẩu mới.

\textbf{Thay đổi 3 -- Trang quản trị user dành cho admin}: Thêm trang
\texttt{users.html} và các endpoint admin tương ứng (\texttt{/waf-api/auth/users/*})
cho phép admin thực hiện đầy đủ CRUD trên tài khoản: tạo mới với role tùy
chọn, sửa email/role, reset password (không cần old\_password do là admin
action), và xóa.

\textbf{Thay đổi 4 -- Last-admin protection}: Trước khi xóa hoặc demote
admin, hệ thống truy vấn \texttt{SELECT COUNT(*) FROM users WHERE role =
'admin'}. Nếu count = 1 và target là chính admin đó, thao tác bị từ chối
với mã 400 và message rõ ràng (\texttt{"cannot delete the last admin"}).
Bất biến này đảm bảo hệ thống luôn có ít nhất một admin -- một tính chất
gọi là \textit{system invariant} trong CS lý thuyết.

\textbf{Thay đổi 5 -- Self-action protection}: Admin không được xóa hoặc
demote chính tài khoản đang đăng nhập. Kiểm tra này được thực hiện ở
\textbf{cả hai tầng}: handler so sánh \texttt{actorID} (từ JWT) với
\texttt{targetID} (trên path), và UI disable nút Delete cho hàng tương ứng
với chính admin đang đăng nhập. Hai tầng phòng thủ thực hiện cùng kiểm
tra -- nếu một tầng có lỗi, tầng còn lại vẫn bắt được, theo nguyên tắc
defense-in-depth.

\textbf{Bổ sung -- JWT trong HttpOnly cookie}: Chuyển từ lưu JWT trong
\texttt{localStorage} (truy cập được từ JavaScript) sang HttpOnly cookie
với \texttt{Secure} flag. JavaScript không đọc được cookie này, ngăn chặn
trộm token qua XSS. Đồng thời vẫn hỗ trợ \texttt{Authorization: Bearer}
header cho API client.

\textbf{Bổ sung -- Audit logging mọi thao tác}: Toàn bộ thao tác auth
(login thành công/thất bại, đổi password, tạo/sửa/xóa user) đều ghi vào
\texttt{audit\_logs} với \texttt{event\_type = 'system'} và
\texttt{actor\_id}. Điều này tạo dấu vết kiểm toán đầy đủ cho mọi can
thiệp vào lớp người dùng.

\subsection{Kết quả}
\label{subsection:5.5.3}

Toàn bộ các thay đổi được kiểm chứng bằng 10 test case auth (bảng
\ref{table:auth-tests}, mục \ref{subsection:4.6.5}) và \textbf{tất cả đều
pass}. Đặc biệt quan trọng:

\begin{itemize}[leftmargin=*]
    \item Test regression \textbf{privilege escalation}: \texttt{POST /register
    \{"role":"admin"\}} trả về user mới có \texttt{role = "viewer"}. Lỗ
    hổng nghiêm trọng nhất đã được đóng vĩnh viễn ở tầng struct.

    \item Bốn test \textbf{invariant lockout}: Xóa/demote admin cuối cùng,
    và tự xóa/demote chính mình -- cả bốn trả về 400 với thông điệp rõ
    ràng. Hệ thống không thể bị đẩy vào trạng thái không có admin.

    \item Test \textbf{change password}: Sai old\_password \(\to\) 401;
    đúng old\_password \(\to\) 200 + audit log entry. Cơ chế chống takeover
    khi JWT bị lộ hoạt động đúng.

    \item Endpoint matrix kiểm tra trên 100\% các endpoint admin: không
    có token \(\to\) 401, viewer/editor token \(\to\) 403, admin token
    \(\to\) 200. Không có endpoint admin nào bỏ sót gate \texttt{requireAdmin}.
\end{itemize}

Ý nghĩa rộng hơn: hệ thống đã chuyển từ trạng thái \textit{có lỗ hổng
critical đã biết} sang trạng thái \textit{tất cả các vector tấn công đã
biết đều được vá và có test regression bảo vệ}. Đây là sự khác biệt giữa
một codebase ``demo'' và một codebase đã sẵn sàng cho môi trường production
sơ cấp.

\bigskip
\noindent\textbf{Kết chương.}
Chương này đã trình bày năm đóng góp độc lập của đề tài: cơ chế anomaly
scoring hai lớp với ML gray-zone activation tận dụng được điểm mạnh của
cả rule lẫn ML mà không phải chịu chi phí latency của ML trên toàn bộ
traffic; bộ luật 78 rule với transform chain JSON-declarative và TOKEN
proximity match chống evasion một cách hệ thống; pipeline huấn luyện với
hard gate đã thực sự ngăn được một regression (v6) khỏi vào production;
thiết kế interface tách biệt engine và ML client mở đường cho test isolation
và khả năng thay thế backend; và đợt củng cố lớp auth đã đóng lỗ leo thang
đặc quyền nghiêm trọng cùng các bất biến bảo vệ chống lockout. Mỗi đóng
góp đều có giá trị độc lập có thể tham khảo trong các dự án WAF mã nguồn
mở khác hoặc các hệ thống bảo mật ứng dụng web tương tự. Chương cuối
(\ref{section:6.1}) sẽ tổng kết mức độ hoàn thành mục tiêu, nhìn nhận các
hạn chế còn tồn tại, và đề xuất hướng phát triển tiếp theo cho hệ thống.

\end{document}
